\chapter{System Architecture}

The network is hierarchical. Local clusters with baselines of tens of kilometres serve regional illuminators and low-altitude scenarios. Regional clusters extend coverage, while long baselines may support HF, high-altitude and space-object observations when common visibility and a compatible illuminator exist.

Each station contains a radome platform, multiband RF chains, local processing, timing and calibration, optical transport, power conditioning, thermal management and a supervisory interface. The central fusion layer associates detections by waveform, time, frequency, Doppler, polarization and geometric consistency. It must distinguish direct emitter observations, bistatic reflections, interference and tracking hypotheses.

\begin{figure}[H]
\centering
\includegraphics[width=0.95\textwidth]{en/fig01_arquitetura_rede.png}
\caption{Distributed sensing architecture: the upper polygon is an illuminator or emitter, the three domes are independent receiving nodes, green paths carry calibrated observations, and the lower purple block is central fusion.}
\label{fig:fig01_arquitetura_rede}
\end{figure}
